\section{Publikationsformen}
%\subsection{Wie gebe ich was in Zotero ein?}
\subsection{Selbstständige Publikationen}
%### Bücher
\subsubsection{Bücher}
Eintragsart: \textit{Buch}.
\pubtable{
\textit{Autor}&Vor- und Nachname sollten getrennt eingegeben werden. Mit dem \verb|+|-Knopf rechts können Sie mehrere Autoren angeben.\\
\textit{Titel}&Achten Sie darauf, Titel und Untertitel mit einem Punkt zu trennen und englische Titel im Title Case anzugeben. Sie können unter dem \verb|...|-Knopf Titel automatisch in Title Case umformen.\\
\textit{Ort}\\
\textit{Datum}&Tragen Sie hier das Jahr ein\\
\textit{Auflage}&Geben Sie hier lediglich eine Zahl ein, so wird diese als hochgestellte Zahl vor dem Jahr ausgegeben. Bei weiteren Informationen (z. B. "`zweite und verbesserte Auflage"') werden diese als eigenständiger Text ausgegeben.
}
\subsubsection{Bücher in gezählten Reihen}
Um eine gezählte Reihe anzugeben, nutzen Sie die Felder \textit{Reihe} und \textit{Nummer der Reihe}. Die Reihe wird nur dann ausgegeben, wenn beide Felder gesetzt wurden. 
\subsubsection{Unveröffentlichte oder online publizierte Hochschulschriften}
Eintragsart: \textit{Dissertation}.
\pubtable{
\textit{Art}&Hier können Sie z. B. "`Diss"' oder "`Masterarbeit"' angeben\\
\textit{Universität}&Geben Sie hier die Hochschule oder Universität ein, bei der Arbeit eingereicht wurde.
}
\subsubsection{Reprints}
Reprints lassen sich momentan in Zotero nicht vollständig abbilden. Als Workaround können Sie das Feld \textit{Auflage} nutzen, um Informationen zum Reprint einzugeben, z. B. "`Faksimile-Nachdruck unter dem Titel Musikalisches Lexikon"'.
\subsection{Unselbstständige Publikationen}
%\subsubsection{Lexikonartikel}
Eintragsart: \textit{Enzyklopädieartikel}.
\pubtable{
\textit{Titel}&Geben Sie hier den Titel des Artikels an \ldots\\
\textit{Titel der Enzyklopädie}&\ldots und hier den Titel des Lexikons bzw. der Enzyklopädie.\\
\textit{Autor} und \textit{Herausgeber}&Direkt unter dem Titel können Sie den Namen des Autors angeben. Um den Herausgeber anzugeben, klicken Sie auf den \verb|+|-Knopf rechts. Nun sollte ein weiteres Feld "`Herausgeber"' erscheinen.
}
\subsubsection{Hinweise zu speziellen Lexika}
\subsubsection*{MGG2}
\pubtable{
\textit{Auflage}&"`2., neubearbeitete Ausgabe"'\\
\textit{Band}&Tragen Sie hier auch den Teil der MGG ein, also z. B. "`Sachteil, Bd. 5"'\\
\textit{Seiten}&Geben Sie hier die Spaltenzahlen mit dem Zusatz "`Sp."' ein. Zotero erkennt, dass es sich um eine nicht-numerische Seitenzahl handelt und wird den ansonsten automatischen Zusatz "`S. [\ldots]"' weglassen.
}
\subsubsection*{NG2}
\pubtable{
\textit{Auflage}&"`2., neubearbeitete Ausgabe"' eintragen.
}
\subsubsection{Artikel in Onlinelexika (GMO)}
Lassen Sie die Felder \textit{Ort} und \textit{Datum} leer und nutzen Sie stattdessen: \textit{Titel}, \textit{Autor}, \textit{Titel der Enzyklopädie}, \textit{URL} und \textit{Heruntergeladen am}.
\subsubsection{Aufsätze in Sammelpublikationen}
Nutzen Sie \textit{Buchteil} als Eintragsart. 
\textit{Aufsätze in Kongressberichten}
Nutzen Sie \textit{Konferenzpaper} als Eintragsart. Informationen zu Titel und Ort des Kongresses nehmen in den Titel des Konferenzbandes auf, \textbf{nicht} in die separat dafür vorgesehenen Felder.
\textit{Aufsätze in Festschriften}
Nutzen Sie die normale Form (s. o. "`Aufsätze in Sammelpublikationen"').
\textit{Aufsätze in wissenschaftlichen Zeitschriften oder Jahrbüchern}
Nutzen Sie "`Zeitschriftenartikel"' als Eintragsart.
\pubtable{
\textit{Publikation}&Hier tragen Sie den Namen der Zeitschrift, \ldots\\
\textit{Band}&\ldots hier den Jahrgang \ldots\\
\textit{Ausgabe}&\ldots und hier das Heft.\\
\textit{Zeitschriftenabkürzung}&Es ist bei bekannten und oft zitierten Zeitschriften sinnvoll, dieses Feld zu setzen. Eine Übersicht über geläufige Abkürzungen finden Sie im ersten Band der MGG2.\\
\textit{DOI}&Bei Onlinezeitschriften sollten Sie hier die DOI des Aufsatzes eintragen, also z. B. \verb|10.31751/1188|. Beim Zitieren wird automatisch ein vollständiger Link ergänzt.\\
\textit{URL}&Sofern eine DOI gesetzt ist, wird dieses Feld beim Zitieren ignoriert. Es ist nur relevant, wenn keine DOI verknüpft ist.
}
\subsubsection{Artikel in Zeitungen oder Magazinen}
Eintragsart: \textit{Zeitungsartikel} (Achtung, nicht Zeitschriftenartikel).
\pubtable{
\textit{Publikation}&Name der Zeitung.
}
\newpage % TODO: care 
\subsubsection{Dokumente in Briefausgaben oder Dokumentensammlungen}
Nehmen Sie hier nur die Briefausgabe oder Dokumentensammlung selbst in Zotero auf (als \textit{Buch}). Beim Zitieren geben Sie manuell die Angaben zum entsprechenden Dokument an und nutzen Zotero lediglich für den zweiten Teil der Angabe, z. B. 
\begin{displayquote}
Brief von Unger an Graf Esterhazy, Wien, 20. August 1820, zitiert nach: [hier Briefausgabe hier mit Zotero einfügen].
\end{displayquote}
%
